% Options for packages loaded elsewhere
\PassOptionsToPackage{unicode}{hyperref}
\PassOptionsToPackage{hyphens}{url}
\documentclass[
]{article}
\usepackage{xcolor}
\usepackage[margin=1in]{geometry}
\usepackage{amsmath,amssymb}
\setcounter{secnumdepth}{-\maxdimen} % remove section numbering
\usepackage{iftex}
\ifPDFTeX
  \usepackage[T1]{fontenc}
  \usepackage[utf8]{inputenc}
  \usepackage{textcomp} % provide euro and other symbols
\else % if luatex or xetex
  \usepackage{unicode-math} % this also loads fontspec
  \defaultfontfeatures{Scale=MatchLowercase}
  \defaultfontfeatures[\rmfamily]{Ligatures=TeX,Scale=1}
\fi
\usepackage{lmodern}
\ifPDFTeX\else
  % xetex/luatex font selection
\fi
% Use upquote if available, for straight quotes in verbatim environments
\IfFileExists{upquote.sty}{\usepackage{upquote}}{}
\IfFileExists{microtype.sty}{% use microtype if available
  \usepackage[]{microtype}
  \UseMicrotypeSet[protrusion]{basicmath} % disable protrusion for tt fonts
}{}
\makeatletter
\@ifundefined{KOMAClassName}{% if non-KOMA class
  \IfFileExists{parskip.sty}{%
    \usepackage{parskip}
  }{% else
    \setlength{\parindent}{0pt}
    \setlength{\parskip}{6pt plus 2pt minus 1pt}}
}{% if KOMA class
  \KOMAoptions{parskip=half}}
\makeatother
\usepackage{color}
\usepackage{fancyvrb}
\newcommand{\VerbBar}{|}
\newcommand{\VERB}{\Verb[commandchars=\\\{\}]}
\DefineVerbatimEnvironment{Highlighting}{Verbatim}{commandchars=\\\{\}}
% Add ',fontsize=\small' for more characters per line
\usepackage{framed}
\definecolor{shadecolor}{RGB}{248,248,248}
\newenvironment{Shaded}{\begin{snugshade}}{\end{snugshade}}
\newcommand{\AlertTok}[1]{\textcolor[rgb]{0.94,0.16,0.16}{#1}}
\newcommand{\AnnotationTok}[1]{\textcolor[rgb]{0.56,0.35,0.01}{\textbf{\textit{#1}}}}
\newcommand{\AttributeTok}[1]{\textcolor[rgb]{0.13,0.29,0.53}{#1}}
\newcommand{\BaseNTok}[1]{\textcolor[rgb]{0.00,0.00,0.81}{#1}}
\newcommand{\BuiltInTok}[1]{#1}
\newcommand{\CharTok}[1]{\textcolor[rgb]{0.31,0.60,0.02}{#1}}
\newcommand{\CommentTok}[1]{\textcolor[rgb]{0.56,0.35,0.01}{\textit{#1}}}
\newcommand{\CommentVarTok}[1]{\textcolor[rgb]{0.56,0.35,0.01}{\textbf{\textit{#1}}}}
\newcommand{\ConstantTok}[1]{\textcolor[rgb]{0.56,0.35,0.01}{#1}}
\newcommand{\ControlFlowTok}[1]{\textcolor[rgb]{0.13,0.29,0.53}{\textbf{#1}}}
\newcommand{\DataTypeTok}[1]{\textcolor[rgb]{0.13,0.29,0.53}{#1}}
\newcommand{\DecValTok}[1]{\textcolor[rgb]{0.00,0.00,0.81}{#1}}
\newcommand{\DocumentationTok}[1]{\textcolor[rgb]{0.56,0.35,0.01}{\textbf{\textit{#1}}}}
\newcommand{\ErrorTok}[1]{\textcolor[rgb]{0.64,0.00,0.00}{\textbf{#1}}}
\newcommand{\ExtensionTok}[1]{#1}
\newcommand{\FloatTok}[1]{\textcolor[rgb]{0.00,0.00,0.81}{#1}}
\newcommand{\FunctionTok}[1]{\textcolor[rgb]{0.13,0.29,0.53}{\textbf{#1}}}
\newcommand{\ImportTok}[1]{#1}
\newcommand{\InformationTok}[1]{\textcolor[rgb]{0.56,0.35,0.01}{\textbf{\textit{#1}}}}
\newcommand{\KeywordTok}[1]{\textcolor[rgb]{0.13,0.29,0.53}{\textbf{#1}}}
\newcommand{\NormalTok}[1]{#1}
\newcommand{\OperatorTok}[1]{\textcolor[rgb]{0.81,0.36,0.00}{\textbf{#1}}}
\newcommand{\OtherTok}[1]{\textcolor[rgb]{0.56,0.35,0.01}{#1}}
\newcommand{\PreprocessorTok}[1]{\textcolor[rgb]{0.56,0.35,0.01}{\textit{#1}}}
\newcommand{\RegionMarkerTok}[1]{#1}
\newcommand{\SpecialCharTok}[1]{\textcolor[rgb]{0.81,0.36,0.00}{\textbf{#1}}}
\newcommand{\SpecialStringTok}[1]{\textcolor[rgb]{0.31,0.60,0.02}{#1}}
\newcommand{\StringTok}[1]{\textcolor[rgb]{0.31,0.60,0.02}{#1}}
\newcommand{\VariableTok}[1]{\textcolor[rgb]{0.00,0.00,0.00}{#1}}
\newcommand{\VerbatimStringTok}[1]{\textcolor[rgb]{0.31,0.60,0.02}{#1}}
\newcommand{\WarningTok}[1]{\textcolor[rgb]{0.56,0.35,0.01}{\textbf{\textit{#1}}}}
\usepackage{graphicx}
\makeatletter
\newsavebox\pandoc@box
\newcommand*\pandocbounded[1]{% scales image to fit in text height/width
  \sbox\pandoc@box{#1}%
  \Gscale@div\@tempa{\textheight}{\dimexpr\ht\pandoc@box+\dp\pandoc@box\relax}%
  \Gscale@div\@tempb{\linewidth}{\wd\pandoc@box}%
  \ifdim\@tempb\p@<\@tempa\p@\let\@tempa\@tempb\fi% select the smaller of both
  \ifdim\@tempa\p@<\p@\scalebox{\@tempa}{\usebox\pandoc@box}%
  \else\usebox{\pandoc@box}%
  \fi%
}
% Set default figure placement to htbp
\def\fps@figure{htbp}
\makeatother
\setlength{\emergencystretch}{3em} % prevent overfull lines
\providecommand{\tightlist}{%
  \setlength{\itemsep}{0pt}\setlength{\parskip}{0pt}}
\usepackage{bookmark}
\IfFileExists{xurl.sty}{\usepackage{xurl}}{} % add URL line breaks if available
\urlstyle{same}
\hypersetup{
  pdftitle={The Savage-Dickey Density Ratio},
  hidelinks,
  pdfcreator={LaTeX via pandoc}}

\title{The Savage-Dickey Density Ratio}
\author{}
\date{\vspace{-2.5em}2026-04-17}

\begin{document}
\maketitle

\subsection{Introduction}\label{introduction}

The Savage-Dickey density ratio is a small but powerful identity: when a
point-null hypothesis is nested inside a continuous alternative, the
Bayes factor in favour of the null can be read off as the ratio of the
marginal posterior to the marginal prior, both evaluated at the null
value. This sidesteps the notoriously hard computation of the marginal
likelihood under each model and turns Bayes-factor calculation into
something you can do from MCMC draws and a one-line evaluation of the
prior density. It is the Bayesian counterpart of asking ``did the data
move probability mass away from the null?'' and quantifying the answer.

\subsection{Prerequisites}\label{prerequisites}

A working understanding of Bayes factors, posterior densities, and the
difference between point-null and interval-null hypotheses. Familiarity
with density estimation from MCMC samples is also useful.

\subsection{Theory}\label{theory}

For a point null \(H_0 : \theta = \theta_0\) nested in a continuous
alternative \(H_1 : \theta \in \Theta\), with the prior under \(H_1\)
taken as the marginal prior \(p(\theta)\) on the parameter of interest,
the Bayes factor in favour of the null is

\[\mathrm{BF}_{01} = \frac{p(\theta_0 \mid y)}{p(\theta_0)}.\]

If the posterior puts more density at the null than the prior did, the
data support \(H_0\); if it puts less, they support \(H_1\). The
quantity is interpretable on Jeffreys' evidence scale and is invariant
under monotone reparameterisations of nuisance parameters because the
alternative's prior must be set up consistently. The identity assumes
the prior under \(H_1\) reduces to the prior under \(H_0\) when
\(\theta = \theta_0\); this is automatic for nested-null setups but can
fail under improper or hierarchical priors.

\subsection{Assumptions}\label{assumptions}

A continuous parameter, the null nested inside the alternative, and a
proper prior whose density at \(\theta_0\) is finite and known. Improper
priors invalidate the construction because the Bayes factor depends on a
normalisation constant that has been thrown away.

\subsection{R Implementation}\label{r-implementation}

\begin{Shaded}
\begin{Highlighting}[]
\FunctionTok{library}\NormalTok{(polspline)}

\FunctionTok{set.seed}\NormalTok{(}\DecValTok{2026}\NormalTok{)}
\CommentTok{\# Posterior samples (e.g., from MCMC)}
\NormalTok{theta\_post }\OtherTok{\textless{}{-}} \FunctionTok{rnorm}\NormalTok{(}\DecValTok{5000}\NormalTok{, }\AttributeTok{mean =} \FloatTok{0.2}\NormalTok{, }\AttributeTok{sd =} \FloatTok{0.3}\NormalTok{)}
\NormalTok{theta\_prior\_density }\OtherTok{\textless{}{-}} \FunctionTok{dnorm}\NormalTok{(}\DecValTok{0}\NormalTok{, }\AttributeTok{mean =} \DecValTok{0}\NormalTok{, }\AttributeTok{sd =} \DecValTok{1}\NormalTok{)   }\CommentTok{\# N(0, 1) prior at 0}

\CommentTok{\# Posterior density at 0 via logspline fit}
\NormalTok{logspl }\OtherTok{\textless{}{-}} \FunctionTok{logspline}\NormalTok{(theta\_post)}
\NormalTok{post\_density\_at\_0 }\OtherTok{\textless{}{-}} \FunctionTok{dlogspline}\NormalTok{(}\DecValTok{0}\NormalTok{, logspl)}

\NormalTok{BF01 }\OtherTok{\textless{}{-}}\NormalTok{ post\_density\_at\_0 }\SpecialCharTok{/}\NormalTok{ theta\_prior\_density}
\NormalTok{BF01; }\DecValTok{1}\SpecialCharTok{/}\NormalTok{BF01}
\end{Highlighting}
\end{Shaded}

\subsection{Output \& Results}\label{output-results}

The ratio of two scalars, ideally accompanied by its inverse so both
directions are visible. A logspline density fit to the posterior draws
is a robust choice: it captures skew and tail behaviour better than a
normal approximation and is less noisy than a kernel density estimate at
the boundary.

\subsection{Interpretation}\label{interpretation}

A reporting sentence: ``The Savage-Dickey Bayes factor was
\(\mathrm{BF}_{10} = 3.1\) under a \(\mathcal N(0, 1)\) prior on the
alternative --- moderate evidence for a non-zero effect --- versus
\(\mathrm{BF}_{01} = 0.32\).'' On Jeffreys' scale,
\(1 < \mathrm{BF} < 3\) is anecdotal, \(3\)--\(10\) is moderate,
\(10\)--\(30\) strong, and beyond \(30\) very strong; the symmetric
thresholds apply to \(\mathrm{BF}_{01}\) for evidence in favour of the
null.

\subsection{Practical Tips}\label{practical-tips}

\begin{itemize}
\tightlist
\item
  Always disclose the prior under the alternative; the Bayes factor is
  sensitive to its scale, particularly via the Lindley-Bartlett paradox
  where a very wide prior favours the null mechanically.
\item
  Use a logspline or kernel density estimate that is robust at the
  boundary --- naive Gaussian approximations break down when the
  posterior is skewed.
\item
  Confirm enough posterior draws fall near \(\theta_0\) for the density
  estimate to be stable; if the data strongly support \(H_1\), the
  posterior density at the null is small and noisy. Increase chain
  length when reporting strong evidence against \(H_0\).
\item
  For interval nulls or composite hypotheses, the Savage-Dickey identity
  does not apply; use bridge sampling
  (\texttt{bridgesampling::bridge\_sampler()}) instead.
\item
  Prefer a sensitivity analysis: report the Bayes factor under at least
  two prior scales, so readers can judge how prior-driven the conclusion
  is.
\end{itemize}

\subsection{R Packages Used}\label{r-packages-used}

\texttt{polspline} for the logspline posterior density estimate at the
null, \texttt{bayestestR} for an integrated Savage-Dickey wrapper that
works directly with \texttt{brms} and \texttt{rstanarm} fits, and
\texttt{bridgesampling} for the more general marginal-likelihood
approach when the Savage-Dickey shortcut does not apply.

\subsection{For Reviewers}\label{for-reviewers}

\textbf{What to look for in a paper using this method.}

\begin{itemize}
\tightlist
\item
  \textbf{Common misapplications.}
\item
  \textbf{Diagnostics that should be reported but often aren't.}
\item
  \textbf{Red flags in tables and figures.}
\item
  \textbf{What to verify.}
\item
  \textbf{What an adequate Methods paragraph must contain.}
\end{itemize}

\subsection{See also --- labs in this
chapter}\label{see-also-labs-in-this-chapter}

\begin{itemize}
\tightlist
\item
  \href{labs/lab-bayesian-thinking-control-vs-decision-error.Rmd}{Bayesian
  thinking; α control vs decision error}
\item
  \href{labs/lab-brms-stan-loo-hierarchical-models.Rmd}{brms / Stan,
  LOO, hierarchical models}
\end{itemize}

\end{document}
